\chapter*{How to Use This Book}
\label{frontmatter:how-to-use}
\addcontentsline{toc}{chapter}{How to Use This Book}

\section*{The reading rhythm}

A serious reader takes between twenty-four and thirty-six months to traverse \emph{Engineering} once, including the projects and the exercises. We have known faster readings to be unrealistic and slower readings to be honest in their own way. The reader who finishes in eighteen months has skipped the projects. The reader who finishes in seventy-two months has built a parallel life around the book. We respect both.

The book offers three pacing tracks: \emph{slow} (thirty-six months at fifteen pages a day), \emph{standard} (twenty-four months at thirty pages a day, the default), and \emph{intensive} (eighteen months at forty-five pages a day, requiring something close to a part-time commitment). We offer no opinion on which the reader should choose; we offer only the schedules.

\section*{The four kinds of unit}

The book contains four kinds of unit. Each is marked in its own environment.

\begin{itemize}
  \item \engterm{Exposition}: prose paragraphs, derivations, worked examples. Read these.
  \item \engterm{Exercises}: short problems closing each chapter. We expect the reader to solve at least half. The exercise solutions are not the point; the act of solving them is.
  \item \engterm{Projects}: longer build-and-measure assignments closing each volume. Build them. The reader who reads the projects without building the projects becomes an informed spectator, not an engineer.
  \item \engterm{Labs}: physical or computational experiments embedded in the prose. Run them. Document the failures. Failed labs are the most valuable kind.
\end{itemize}

We aim for roughly one third of the book to be exposition; the remaining two thirds is exercises, projects, labs, and case studies. The book is what the reader does with it.

\section*{The role of failure}

Volume X treats failure as a domain in its own right. Failure also runs through every other volume, in a chapter-by-chapter section that asks how this model breaks, how this material fatigues, how this system collapses, and how the engineer notices in time. The reader will get derivations wrong. The reader will assemble circuits that smoke. The reader will write simulators whose answers diverge. The reader will measure the wrong quantity. We treat each of these as a learning event. The reader who treats failure as an interruption of learning has misunderstood what learning is.

A laboratory notebook accompanies the book. We expect the reader to keep one. Failures are recorded with date, condition, what was expected, what was observed, and the correction. This habit is the single highest-leverage skill the book transmits, and the only one no chapter is dedicated to.

\section*{Reading non-sequentially}

The book can be read non-sequentially with a price.

\begin{itemize}
  \item Volumes I and II are prerequisites for all subsequent volumes. We do not allow them to be skipped.
  \item Volume III is prerequisite for Volumes IV, VIII, and IX.
  \item Volume IV is prerequisite for Volumes V, VIII, and XII.
  \item Volume V is prerequisite for Volumes VIII and XII.
  \item Volume VI is independent of VII; either may be read before the other after I-IV.
  \item Volume VII is prerequisite for Volumes VIII (electronics) and IX (control).
  \item Volumes VIII, IX, X are read in order, after I-VII.
  \item Volumes XI and XII are read last.
\end{itemize}

A reader with prior coursework in some domain may skim its book and concentrate on its exercises. We expect this to be done honestly. Self-deception about what one already knows is the most expensive form of skipping.

\section*{Tools the reader needs}

A laboratory notebook (paper). A calculator that can do floating-point arithmetic to at least eight digits. A computer with a UNIX-like environment. Python with NumPy, SciPy, Matplotlib, SymPy. Access to a workshop with a vise, a soldering iron, a multimeter, a 3D printer, basic hand tools. A camera for documentation. A network connection for primary sources.

We list, in each chapter's project, the specific additional tools required. None is exotic; none is expensive at the scale a serious reader can afford. The required investment in tools is small. The required investment in time is the entire price.

The reader who lacks a workshop should follow the simulation/household track that runs in parallel to the physical track. The simulation track preserves the discipline of the projects without requiring soldering equipment, milling machines, or printed-circuit-board fabrication. We mark each project with a physical or simulation tag and offer alternatives where possible.

\section*{When to put the book down}

The book is dense. The reader will reach passages where the meaning slips. We recommend putting the book down, walking away, returning the next morning, and rereading. The reader who insists on understanding every paragraph on first reading reads slower and learns less. The reader who rereads on a day-scale cadence learns more, faster.

The reader who, after three honest re-readings, still does not understand a paragraph, marks the paragraph and moves on. Most such paragraphs become clear in subsequent chapters. The few that do not are recorded in the laboratory notebook as unresolved questions. Engineering is built on accumulated unresolved questions, not on the absence of them.
